\documentclass[12pt]{article}
\usepackage{amsmath, amssymb}
\usepackage{geometry}
\geometry{margin=1in}

\title{CSE400: Fundamentals of Probability in Computing\\
Lecture 8: Gaussian, Uniform, Exponential, Laplace, and Gamma Random Variables}
\author{}
\date{}

\begin{document}
\maketitle

\section*{Gaussian Random Variable}

\subsection*{Definition}

A continuous random variable $X$ is said to be a \textbf{Gaussian (Normal) random variable} if its probability density function (PDF) can be written as
\[
f_X(x) = \frac{1}{\sqrt{2\pi\sigma^2}} \exp\left(-\frac{(x-m)^2}{2\sigma^2}\right),
\]
where
\[
m = \mathbb{E}[X], \qquad \sigma^2 = \text{Var}(X).
\]

Notation:
\[
X \sim \mathcal{N}(m,\sigma^2).
\]

\subsection*{Properties}

\begin{itemize}
\item PDF is symmetric about $x=m$.
\item Mean equals median equals mode.
\item Total area under the PDF equals 1.
\end{itemize}

\section*{Standard Forms}

\subsection*{Error Function}

\[
\mathrm{erf}(x) = \frac{2}{\sqrt{\pi}} \int_{0}^{x} e^{-t^2} dt
\]

\subsection*{Complementary Error Function}

\[
\mathrm{erfc}(x) = 1 - \mathrm{erf}(x)
= \frac{2}{\sqrt{\pi}} \int_{x}^{\infty} e^{-t^2} dt
\]

\subsection*{$\Phi$-Function (CDF of Standard Normal)}

\[
\Phi(x) = \frac{1}{\sqrt{2\pi}} \int_{-\infty}^{x} e^{-t^2/2} dt
\]

\subsection*{$Q$-Function (Gaussian Tail Function)}

\[
Q(x) = \frac{1}{\sqrt{2\pi}} \int_{x}^{\infty} e^{-t^2/2} dt
\]

\subsection*{Identity}

\[
Q(x) = 1 - \Phi(x)
\]

\section*{Relation Between $\Phi$-Function and $Q$-Function}

\subsection*{Evaluating the CDF}

\[
F_X(x) = \int_{-\infty}^{x} \frac{1}{\sqrt{2\pi\sigma^2}}
\exp\left(-\frac{(y-m)^2}{2\sigma^2}\right) dy
\]

Let
\[
t = \frac{y-m}{\sigma}
\]

Then
\[
F_X(x) = \int_{-\infty}^{\frac{x-m}{\sigma}}
\frac{1}{\sqrt{2\pi}} e^{-t^2/2} dt
= \Phi\left(\frac{x-m}{\sigma}\right)
\]

\subsection*{Tail Probability}

\[
\Pr(X > x) =
\int_{\frac{x-m}{\sigma}}^{\infty}
\frac{1}{\sqrt{2\pi}} e^{-t^2/2} dt
= Q\left(\frac{x-m}{\sigma}\right)
\]

\subsection*{Key Identities}

\[
Q(x) = 1 - \Phi(x)
\]
\[
F_X(x) = 1 - Q\left(\frac{x-m}{\sigma}\right)
\]

\section*{Expectation of a Function of a Random Variable}

\subsection*{Continuous Random Variable}

\[
\mathbb{E}[g(X)] = \int g(x) f_X(x) \, dx
\]

\subsection*{Discrete Random Variable}

\[
\mathbb{E}[g(X)] = \sum_k g(x_k) P_X(x_k)
\]

\subsection*{Linearity of Expectation}

For constants $a,b$,
\[
\mathbb{E}[aX + b] = a\mathbb{E}[X] + b
\]

For random variables $X_1,\dots,X_N$,
\[
\mathbb{E}\left[\sum_{k=1}^{N} g_k(X_k)\right]
= \sum_{k=1}^{N} \mathbb{E}[g_k(X_k)]
\]

\section*{Higher-Order Moments}

\subsection*{Second Central Moment}

\[
\sigma^2 = \mathbb{E}[(X-\mu)^2]
\]

\subsection*{Skewness}

\[
C_s = \frac{\mathbb{E}[(X-\mu)^3]}{\sigma^3}
\]

Measures symmetry about the mean.

\subsection*{Kurtosis}

\[
C_k = \frac{\mathbb{E}[(X-\mu)^4]}{\sigma^4}
\]

Large value indicates sharper peak near the mean.

\section*{Example: Gaussian Random Variable}

A random variable $X$ has PDF

\[
f_X(x) = \frac{1}{\sqrt{8\pi}}
\exp\left(-\frac{(x+3)^2}{8}\right)
\]

\subsection*{Identification}

Comparing with

\[
f_X(x) = \frac{1}{\sqrt{2\pi\sigma^2}}
\exp\left(-\frac{(x-m)^2}{2\sigma^2}\right)
\]

we identify

\[
m=-3, \qquad \sigma^2=4, \qquad \sigma=2
\]

\subsection*{Gaussian Relations}

\[
F_X(x) = \Phi\left(\frac{x-m}{\sigma}\right)
\]
\[
\Pr(X>x)=Q\left(\frac{x-m}{\sigma}\right)
\]

\subsection*{1) $\Pr(X\le 0)$}

\[
\Pr(X\le0)=F_X(0)
=\Phi\left(\frac{0-(-3)}{2}\right)
=\Phi\left(\frac{3}{2}\right)
=1-Q\left(\frac{3}{2}\right)
\]

\subsection*{2) $\Pr(X>4)$}

\[
\Pr(X>4)
=Q\left(\frac{4-(-3)}{2}\right)
=Q\left(\frac{7}{2}\right)
\]

\end{document}
